\begin{table}[!htbp]
\centering
\caption{Computer-use and broad-family conditioning on common support}\label{tab:characteristics}
\resizebox{\textwidth}{!}{%
\begin{tabular}{llrrr}
\toprule
Family-month paths & Computer-use control & Occupations & Q5 by post [95\% CI] & Change from pooled baseline [95\% CI] \\
\midrule
No & No & 455 & $-0.0958\;[-0.1807,-0.0108]$ & --- \\
No & Yes & 455 & $-0.1966\;[-0.3161,-0.0771]$ & $-0.1009\;[-0.1772,-0.0245]$ \\
Yes & No & 455 & $\phantom{-}0.0354\;[-0.0859,0.1567]$ & $\phantom{-}0.1311\;[0.0352,0.2271]$ \\
Yes & Yes & 455 & $-0.0467\;[-0.1821,0.0888]$ & $\phantom{-}0.0491\;[-0.0642,0.1624]$ \\
\bottomrule
\end{tabular}}
\tabnote{All four post-outcome exploratory rows use the same 455 occupations, rebuilt exposure labels, outcome cells, calendar, and 9,999 common occupation multipliers.  The computer-use rating is standardized with preperiod employment weights and interacted with young $\times$ post.  Family-month rows give every SOC2 family its own young-relative calendar-month path.  The final column uses covariance-preserving paired draws.  The four projections are not additive components, and no residual coefficient is a causal AI effect.}
\end{table}
